\chapter{Thiết kế hệ đa tác tử}\label{chap:1}

Chương này tóm lược tài liệu hướng dẫn của học phần \emph{Developing Multiagent Systems: A Practical Tutorial
with Python and PyTorch}~\cite{ch1:tutorial}, bổ sung bằng hai giáo trình kinh điển về tác tử~\cite{ch1:wooldridge,ch1:russell}.
Mục tiêu là xác lập bộ khái niệm (tác tử, môi trường, PEAS, vai trò, thông điệp, điều phối, kiểm thử) mà
Chương~\ref{chap:4} dùng để mô hình hóa thị trường hai cửa hàng và Chương~\ref{chap:8} dùng để cài đặt hệ mua sắm.

\section{Tác tử, môi trường và mô tả PEAS}

\begin{dinhnghia}{Định nghĩa 1.1 (Tác tử và hệ đa tác tử)}
\textbf{Tác tử} (agent) là một thành phần tính toán nhận thông tin từ môi trường qua cảm biến (sensors), chọn
hành động theo mục tiêu và tác động trở lại môi trường qua cơ cấu chấp hành (actuators). \textbf{Hệ đa tác tử}
(multi-agent system, MAS) gồm từ hai tác tử trở lên có quyết định ảnh hưởng lẫn nhau, qua thông điệp, qua trạng
thái dùng chung hoặc qua chính môi trường~\cite{ch1:wooldridge}.
\end{dinhnghia}

Điểm quan trọng của định nghĩa là các tác tử \emph{không nhất thiết trò chuyện} với nhau. Hai cửa hàng trực tuyến
không gửi thông điệp nào cho nhau, nhưng quyết định giảm giá của cửa hàng này làm thay đổi xác suất mua hàng ở
cửa hàng kia; tương tác gián tiếp qua môi trường như vậy vẫn tạo thành một MAS.

Để đặc tả một tác tử trước khi lập trình, tutorial dùng khung \textbf{PEAS}~\cite{ch1:russell}:
\emph{Performance} (tiêu chí đánh giá), \emph{Environment} (môi trường hoạt động), \emph{Actuators} (cách tác động)
và \emph{Sensors} (thông tin nhận được). Hình~\ref{fig:ch1-peas} đặt PEAS lên vòng cảm nhận--hành động của tác tử
bán hàng trong bài thí nghiệm.

\hinhtikz{ch1_peas}{Vòng cảm nhận -- quyết định -- hành động của tác tử bán hàng, đánh dấu bốn thành phần PEAS.}{ch1-peas}

\nhanxet{Cùng một vòng lặp áp dụng cho cả hai tác tử của bài thí nghiệm; chúng chỉ khác ở khối ``chính sách'':
tác tử học tăng cường cập nhật chính sách từ phần thưởng, tác tử theo luật giữ nguyên một bảng if--else.}

\begin{table}[!htbp]\centering\caption{Đặc tả PEAS của tác tử bán hàng trong bài thí nghiệm giảm giá.}\label{tab:ch1-peas}
{\small\begin{tabular}{@{}L{3.2cm}L{11.3cm}@{}}\toprule
Thành phần & Nội dung \\\midrule
Performance & Lợi nhuận trung bình mỗi bước trong pha đánh giá; phụ: tỷ lệ giảm giá, độ ổn định qua các seed \\
Environment & Thị trường mô phỏng: nhu cầu ngẫu nhiên hai mức, một đối thủ, số khách và doanh số ngẫu nhiên \\
Actuators   & Hai hành động: giữ giá hoặc giảm giá trong bước hiện tại \\
Sensors     & Mức nhu cầu hiện tại và hành động của đối thủ ở bước trước \\\bottomrule
\end{tabular}}
\end{table}

Ngoài PEAS, tutorial nêu năm \textbf{thuộc tính} thường gặp của tác tử: tự chủ (autonomy), phản ứng (reactivity),
chủ động (pro-activeness), giao tiếp (social ability) và học (learning). Các thuộc tính này xuất hiện ở mức độ khác
nhau: tác tử theo luật vẫn tự chủ và phản ứng với trạng thái mới nhưng không học; tác tử Q-learning học trong pha
huấn luyện rồi chỉ thực thi chính sách đã học trong pha đánh giá. Sự tách hai pha này là nguyên tắc xuyên suốt của
Chương~\ref{chap:5}.

\section{Phân chia vai trò và hợp đồng thông điệp}

Một MAS được thiết kế tốt bắt đầu bằng việc chia trách nhiệm thành các \textbf{vai trò} (roles) hẹp, mỗi vai trò
có đầu vào và đầu ra xác định. Hệ mua sắm của tutorial gồm bốn vai trò~\cite{ch1:tutorial}:
\textbf{CatalogAgent} tìm sản phẩm ứng viên theo truy vấn; \textbf{RecommendationAgent} lọc theo ngân sách và xếp
hạng; \textbf{InventoryAgent} xác minh tồn kho; \textbf{Coordinator} nhận yêu cầu, gọi lần lượt các vai trò và
kiểm tra kết quả trước khi trả lời người dùng (Hình~\ref{fig:ch1-shop}).

\hinhtikz{ch1_shop}{Hệ mua sắm bốn vai trò của tutorial: mọi lời gọi đi qua Coordinator; mỗi mũi tên là một thông điệp có kiểu.}{ch1-shop}

\nhanxet{Coordinator là điểm kiểm soát duy nhất: nó kiểm tra người nhận, kiểu phản hồi, ngân sách và độ mới của
dữ liệu tồn kho, nên một tác tử trả lời sai bị chặn trước khi lan sang bước sau.}

Các tác tử trao đổi qua \textbf{thông điệp} có cấu trúc cố định. Tutorial dùng bốn trường bắt buộc:
\texttt{sender}, \texttt{receiver}, \texttt{kind} (loại thông điệp, ví dụ \texttt{search\_request},
\texttt{rank\_response}) và \texttt{payload} (dữ liệu). Hệ lớn hơn bổ sung mã tác vụ, mã thông điệp, thời điểm gửi
và hạn chót để truy vết. Cố định hợp đồng thông điệp có hai lợi ích: mỗi tác tử kiểm thử được độc lập bằng thông
điệp giả, và Coordinator phát hiện được phản hồi sai kiểu thay vì lan truyền lỗi. Chương~\ref{chap:8} viết hợp đồng
này thành kiểu dữ liệu (\texttt{dataclass}) được kiểm tra khi nhận.

\section{Điều phối tập trung và phân tán}

Cơ chế \textbf{điều phối} (coordination) quyết định ai gọi ai, theo thứ tự nào và ai giải quyết xung đột.
Điều phối tập trung dễ theo dõi, dễ kiểm thử, nhưng nút điều phối có thể thành điểm nghẽn và điểm lỗi đơn. Điều phối
phân tán tăng quyền tự chủ và khả năng mở rộng, đổi lại phải xử lý xung đột, tính nhất quán và lỗi giao tiếp
(Hình~\ref{fig:ch1-coord}). Hai cửa hàng trong bài thí nghiệm thuộc dạng thứ ba: không có điều phối viên, mỗi tác tử
quyết định độc lập và đồng thời, môi trường áp dụng cả hai hành động rồi trả phần thưởng riêng.

\hinhtikz{ch1_coord}{Ba dạng điều phối: tập trung (hệ mua sắm), phân tán ngang hàng, và tương tác gián tiếp qua môi trường (thị trường hai cửa hàng).}{ch1-coord}

\nhanxet{Dạng điều phối quyết định bài toán học: ở dạng thứ ba, mỗi tác tử thấy đối thủ như một phần của môi
trường; nếu đối thủ cũng học, môi trường trở nên phi dừng (Chương~\ref{chap:2}, Mục~\ref{sec:ch2-challenge}).}

\section{Tích hợp mô hình học và kiểm thử hệ thống}

Mạng nơ-ron trong MAS chỉ là một thành phần dự đoán nhãn (ví dụ ý định của khách) hoặc điểm xếp hạng; tutorial
minh họa bằng mạng \texttt{IntentNet} (PyTorch) và yêu cầu đánh giá nó riêng trên tập kiểm tra~\cite{ch1:tutorial}.

Kiểm thử MAS vì vậy có ba tầng: (i) \emph{kiểm thử đơn vị} từng tác tử với thông điệp giả; (ii) \emph{kiểm thử
tích hợp} luồng qua Coordinator, gồm cả \emph{mô phỏng lỗi} (tác tử im lặng, trả sai kiểu, dữ liệu tồn kho cũ);
(iii) \emph{đánh giá hệ thống} bằng chỉ số theo mục tiêu: tỷ lệ hoàn thành tác vụ, số lần vi phạm ngân sách, số
xác nhận tồn kho chưa xác minh, độ trễ và số thông điệp. Với bài thí nghiệm giảm giá, chỉ số hệ thống là lợi nhuận
đo trên pha đánh giá tách biệt, nhiều seed; Bảng~\ref{tab:ch3-map} ở Chương~\ref{chap:3} so sánh đầy đủ hai hệ.
